%! Representing a Categorical Variable — Unit 1, Lesson 1.2 — Homework (Answer Key)
% GRADUAL RELEASE: this is the lesson's SCORED practice and the LAST component in the packet,
% started in the closing minutes of the period and finished at home. The AP Practice component
% that precedes it stays assigned-but-unscored. When DeltaMath carries this topic instead, that
% assignment replaces this page and the cover's score cell is struck through.
% \answerspace{H}{} reserves exactly H in the blank; the key passes the answer as the second
% argument, so the two files paginate identically by construction. Keep the macro and every
% \boxguard byte-identical in the blank and the key.%
% GENERATED FROM homework/main.tex — this file is the blank with answers filled in, so the
% \answerspace heights and every \boxguard are identical to it.
% NO teachernote here — teacher prose lives in the lesson plan.
\documentclass[12pt]{article}
\usepackage{apstats-article}
\usepackage{apstats-key}   % pulls in -boxes; do NOT also load -boxes
\usepackage{grading}
\gradingassignment{apstats-u01-l02-hw}
\gradingcourse{advanced_statistics}
\newcommand{\answerspace}[2]{\par\nopagebreak\noindent\begin{minipage}[t][#1][t]{\linewidth}%
  \color{keyred}\bfseries #2\end{minipage}\par}

\begin{document}

\pageheader{Unit 1, Lesson 1.2}{Homework --- Answer Key}
{\small\itshape Write each final answer in its box. Only what is inside a box is graded.}\par
% NAMESTRIP: no \namedateperiod here — the name row lives on the cover only.

\vspace{0.15in}

% Authored BYTE-IDENTICALLY in homework/ and homework_key/ — this box is what keeps the
% blank and the key the same number of pages.
\boxguard[8]
\begin{remindbox}
\small
\textbf{This is your graded homework.} It is scored and due next class. You will start it in the
last minutes of the period, so ask about anything that stalls you before you leave, then finish
it on your own.
\end{remindbox}

\vspace{0.12in}

% ── The scored practice ─────────────────────────────────────────────────────
\boxguard[14]
\begin{notesbox}{Counts, Percents, and Fair Comparisons}
Show the arithmetic wherever a question asks for a number, and write every explanation
\emph{in the context of the problem} --- that is what earns credit on the AP exam, and it is the
standard here.
\begin{enumerate}[leftmargin=1.9em, itemsep=6pt]
  \item A club of 50 members chose a service project: Park cleanup 22, Food bank 15,
        Tutoring 13. Park cleanup is \smash{\gradedblank[id=hw1a, type=short, points=1, answer=44, check=numeric, display={44}]{1.8cm}{0.2in}}\% of the club, and the three percents must
        add to \smash{\gradedblank[id=hw1b, type=short, points=1, answer=100, check=numeric, display={100\%}]{1.4cm}{0.2in}}. Is ``most members chose the park cleanup'' a fair thing to say?
        \answerspace{2.15in}{\normalcolor\mdseries\gradedblank[id=hw1c, type=rubric, points=2,
          display={No --- 44\% is the largest group but not a majority.}]{\linewidth}{2.2cm}}
  \item Two hospitals report flu cases: City General had 340 of its 8{,}500 patients, and
        Ridgeview had 60 of its 1{,}000. Which hospital had more cases, and which had the
        higher rate? Show the two rates.
        \answerspace{2.15in}{\normalcolor\mdseries\gradedblank[id=hw2, type=rubric, points=3,
          display={City General had more cases (340 vs.\ 60), but Ridgeview had the
        higher rate: $340/8500 = 4\%$ against $60/1000 = 6\%$.}]{\linewidth}{2.9cm}}
  \item A bar graph of a school's club memberships has bars of height 40, 25, 20, and 15.
        A second graph of the same data has bars of height 0.40, 0.25, 0.20, and 0.15.
        Explain what changed and what did not.
        \answerspace{2.15in}{\normalcolor\mdseries\gradedblank[id=hw3, type=rubric, points=2,
          display={The first plots counts (40+25+20+15 = 100 members), the second the
        same data as proportions. The bars keep identical relative heights --- only the vertical
        scale and its units changed.}]{\linewidth}{2.5cm}}
  \item A survey of 200 people reports these percents: 34\%, 28\%, 22\%, and 11\%. Explain how
        you can tell, without any other information, that something is wrong.
        \answerspace{2.15in}{\normalcolor\mdseries\gradedblank[id=hw4, type=rubric, points=2,
          display={They add to 95\%, not 100\%. Every person must fall in exactly one
        category, so a category or 5\% of the responses is missing.}]{\linewidth}{2.2cm}}
  \tcbbreak
  \item \textbf{AP-style multiple choice.} Two towns report how residents commute. Brookfield: 900 of
        3{,}000 residents bike. Ashton: 240 of 600 residents bike. Which statement is
        \textbf{correct}? Write its letter in the Letter box, then justify it in a sentence.
        \begin{enumerate}[label=\textbf{(\Alph*)}, leftmargin=2.2em, itemsep=2pt]
          \item Brookfield has both more bikers and the higher biking rate.
          \item Ashton has both more bikers and the higher biking rate.
          \item \ans{Brookfield has more bikers, but Ashton has the higher biking rate.}
          \item Ashton has more bikers, but Brookfield has the higher biking rate.
          \item The two towns cannot be compared, because they are different sizes.
                \hfill\gradedright{\textit{Letter:}~\gradedblank[id=hw5, type=choice, points=1, answer=C]{0.45in}{0.24in}}
        \end{enumerate}
        \answerspace{2.15in}{\normalcolor\mdseries\gradedblank[id=hw5j, type=rubric, points=1,
          display={(C). Brookfield: $900/3000 = 30\%$; Ashton: $240/600 = 40\%$.}]{\linewidth}{2.2cm}}
\end{enumerate}
\end{notesbox}

\vspace{0.12in}

% The next-lesson preview closes the packet, so it lives here rather than in AP Practice.
\boxguard[10]
\begin{spiralbox}
\small
\textbf{Next lesson: Representing a Quantitative Variable with Graphs.} Counting worked here
because every value was one of a few category names. Next time the values are numbers that
barely repeat --- you cannot count how many students scored exactly 73.4 --- so we need a
different set of displays, and the one you choose changes what you see.
\end{spiralbox}

\end{document}
